%==============================================================================
% Sjabloon onderzoeksvoorstel bachproef
%==============================================================================
% Gebaseerd op document class `hogent-article'
% zie <https://github.com/HoGentTIN/latex-hogent-article>

\documentclass{hogent-article}

% Invoegen bibliografiebestand
\addbibresource{voorstel.bib}

% Informatie over de opleiding, het vak en soort opdracht
\studyprogramme{Professionele bachelor toegepaste informatica}
\course{Bachelorproef}
\assignmenttype{Onderzoeksvoorstel}

\academicyear{2025-2026}

\title{Ontwerp van een geautomatiseerd systeem voor security monitoring en incident response in SaaS-omgevingen binnen de DevOps-infrastructuur van Devinity.eu}

\author{Sem Demoen}
\email{sem.demoen@student.hogent.be}

\specialisation{System \& Network Administrator}
%\keywords{Scheme, World Wide Web, $\lambda$-calculus}

\setlength{\parskip}{1em}   % afstand tussen alinea's

\begin{document}
\begin{abstract}
Deze bachelorproef onderzoekt hoe een geautomatiseerd systeem voor security monitoring en incident response kan worden ontworpen om securitydata binnen SaaS-omgevingen effectief te centraliseren, analyseren en rapporteren in een DevOps-infrastructuur.
Het onderzoek vertrekt vanuit een concrete casus bij Devinity.eu, waarbij huidige uitdagingen zoals verspreide tools, vertraagde incidentdetectie en inconsistente rapportering dagelijks voorkomen.
De centrale onderzoeksvraag luidt: ``Hoe kan een geautomatiseerd systeem voor security monitoring en incident response worden ontworpen om securitydata binnen de SaaS-omgeving van Devinity.eu effectief te centraliseren, analyseren en rapporteren in een DevOps-infrastructuur?''
Het doel is een proof-of-concept te ontwikkelen dat logs, events en alerts van SaaS-platforms (zoals Microsoft 365, Google Workspace, Okta) en DevOps-omgevingen verzamelt via API's, webhooks en logforwarders,
deze gegevens normaliseert en analyseert, om vervolgens automatische rapportages te genereren voor de technische teams en het management van het bedrijf.
De methodologie omvat een literatuurstudie, analyse van bestaande security monitoring tools, ontwerp en implementatie van een prototype en evaluatie aan de hand van indicatoren zoals detectiesnelheid, foutpositieven en responstijd.
Het verwachte resultaat is een proof-of-concept dat operationele efficiëntie verhoogt, sneller incidenten detecteert en consistente, bruikbare rapportages oplevert.
De verwachte conclusie is dat een gecentraliseerde en geautomatiseerde aanpak voor security monitoring binnen de SaaS- en DevOps-omgeving haalbaar en effectief is, en een significante verbetering kan betekenen voor de zichtbaarheid, incidentrespons en besluitvorming binnen Devinity.eu.
De meerwaarde voor de doelgroep bestaat uit een framework voor geautomatiseerde security monitoring binnen SaaS en DevOps, dat de organisatie helpt risico's beter te beheersen en hun incident response te optimaliseren. GitHub: \url{https://github.com/SemDemoen/bachproef}
\end{abstract}



\tableofcontents

% De hoofdtekst van het voorstel zit in een apart bestand, zodat het makkelijk
% kan opgenomen worden in de bijlagen van de bachelorproef zelf.
\input{voorstel-inhoud}

\printbibliography[heading=bibintoc]

\end{document}